\chapter{Implementación}

Este capítulo describe las decisiones técnicas tomadas durante la
implementación de FlexiPwn, organizadas según las capas de la arquitectura.

\section{Stack Tecnológico}

% [Pegar aquí el texto del documento stack-tecnologico-flexipwn.pdf
% que ya tienes generado. Incluye: lenguaje Python 3.12, uv,
% decisión por cada capa, tabla resumen al final]

\section{Decisión de Tecnología de Virtualización}

% [Pegar aquí el texto del documento virt-decision-flexipwn.pdf,
% incluyendo la tabla comparativa y la justificación de Docker rootless]

\section{Capa 1: Orquestador de Contenedores}

% [Descripción de DockerRootlessProvider, detección de socket,
% estrategia de baseline con healthcheck, rollback en create()]

\section{Capa 2: Monitoreo Pasivo}

% [Descripción de FilesystemMonitor con container.diff(),
% ProcessMonitor con container.top() y lstart,
% MonitorOrchestrator, y baseline de startup]

\subsection{Limitaciones del Monitoreo por Polling}

El monitor de procesos consulta \texttt{container.top()} cada dos segundos.
Procesos que nacen y terminan dentro de ese intervalo son invisibles para
el sistema. Los vectores no detectables en el estado actual incluyen:
procesos efímeros con duración menor a dos segundos, técnicas de
\textit{PPID spoofing}, y ataques \textit{fileless} mediante llamadas al
sistema como \texttt{ptrace} o \texttt{memfd\_create}. Para el alcance
educativo del proyecto, este nivel de detección es aceptable dado que los
ejercicios no requieren detectar procesos sub-segundo ni técnicas de evasión
avanzadas esperables en un atacante universitario.

De forma análoga, \texttt{container.diff()} retorna una instantánea del
estado acumulado del filesystem desde el inicio del contenedor. Para
distinguir cambios del sistema de inicialización de cambios realizados por
el estudiante, se captura un \textit{baseline} del \textit{diff} una vez
que el contenedor está estable. Si la imagen tiene un \texttt{HEALTHCHECK}
configurado, el sistema espera hasta que el contenedor reporte estado
\texttt{healthy} antes de capturar el baseline; de lo contrario, espera un
delay configurable (por defecto 3 segundos) como mecanismo de respaldo.

\subsection{Monitor de Red}
\label{sec:capa2}

Para escenarios que requieren observar tráfico de red, la Capa 2 instancia 
un cuarto monitor: \texttt{NetworkMonitor}. A diferencia de los tres monitores 
anteriores, este monitor depende de un contenedor sidecar que comparte el 
namespace de red del sistema vulnerable.

\subsubsection*{Contenedor sidecar}
\begin{sloppypar}
    Al crear el entorno, si el escenario define al menos un target de tipo 
    \texttt{network\_*}, \texttt{DockerRootlessProvider} lanza un contenedor 
    adicional basado en la imagen \texttt{nicolaka/netshoot} con la siguiente 
    configuración:
\end{sloppypar}

\begin{itemize}
    \item \textbf{Namespace de red compartido}: mediante 
    \texttt{--network container:\{vulnerable\}}, el sidecar comparte 
    exactamente el mismo namespace de red que el vulnerable. Esto hace 
    visible todo el tráfico del vulnerable desde dentro del sidecar, 
    incluyendo comunicaciones por la interfaz de loopback 
    (\texttt{127.0.0.1}), que de otro modo no serían capturables desde 
    una perspectiva externa.
    
    \item \textbf{Comando de captura}: \texttt{tcpdump -i any -A -n -l}, 
    con un filtro opcional definido en el campo 
    \texttt{environment.capture\_filter} del YAML del escenario 
    (por ejemplo, \texttt{port 3306} para restringir la captura 
    a tráfico MySQL).
    
    \item \textbf{Salida en bind mount}: la captura se escribe en un 
    archivo \texttt{traffic.txt} en un directorio del host con 
    permisos \texttt{0o700}, al que solo tiene acceso el proceso de 
    FlexiPwn.
\end{itemize}

El uso de \texttt{-n} desactiva la resolución de nombres, forzando IPs 
numéricas en el output y evitando que el parser del monitor falle ante 
hostnames como \texttt{localhost}. El flag \texttt{-l} activa el 
line-buffering de tcpdump, reduciendo la latencia entre la escritura 
del paquete y su disponibilidad en el archivo del host.

\subsubsection*{Parsing packet-based}

\texttt{NetworkMonitor} lee el archivo \texttt{traffic.txt} de forma 
incremental en cada ciclo de polling, procesando el output de tcpdump 
como unidades de paquete completo en lugar de línea por línea. Esta 
decisión es necesaria porque el output de \texttt{tcpdump -A} en modo 
\texttt{-i any} mezcla bytes ASCII del payload con bytes binarios de 
los headers TCP/IP, y el carácter \texttt{0x0a} (LF) puede aparecer 
dentro de los headers, partiendo artificialmente el payload en múltiples 
líneas si se procesa con \texttt{splitlines()}.

El monitor detecta las cabeceras de paquete mediante una expresión regular 
que acepta el formato de \texttt{-i any} (que incluye el nombre de interfaz 
y dirección entre el timestamp y la IP) y extrae de ellas las IPs y 
puertos de origen y destino. Las líneas entre cabeceras se acumulan en un 
buffer de payload. Cuando se detecta una nueva cabecera, el buffer del 
paquete anterior se procesa: se extrae todo el texto en rango ASCII 
imprimible (\texttt{0x20--0x7e}), y si el resultado supera los cuatro 
caracteres se emite un \texttt{MonitorEvent} de tipo \texttt{network\_payload} 
con el texto extraído y los metadatos del paquete. Adicionalmente, si la 
cabecera contiene \texttt{Flags [S.]}, se emite un evento 
\texttt{network\_connection} con la IP y puerto de destino.

Una implicación importante del modelo packet-based es que el último 
paquete de una ráfaga de tráfico corta podría quedar atrapado en el 
buffer si no llega una cabecera posterior que dispare su procesamiento. 
Para evitar esto, al final de cada ciclo de polling el monitor procesa 
el paquete en curso si el buffer no está vacío, sin esperar al siguiente 
ciclo.

\section{Capa 3: Motor de Evaluación}

% [Descripción del EvaluationEngine, tipos de target implementados,
% condiciones any/all, evaluadores de Capa 3]

Los tipos de target implementados en el estado actual del sistema son:

\begin{description}
    \item[\texttt{file\_created}, \texttt{file\_modified}, \texttt{file\_exists}] 
    Operan sobre eventos del monitor de filesystem. Permiten detectar 
    creación o modificación de archivos en rutas absolutas o mediante 
    patrones glob.
    
    \item[\texttt{process\_running}] Opera sobre eventos del monitor de 
    procesos. Detecta procesos activos por effective UID y substring 
    en el comando.
    
    \item[\texttt{log\_pattern}] Opera sobre eventos del monitor de logs. 
    Aplica un diccionario de campo→regex sobre los campos de la entrada 
    de log parseada como JSON.
    
    \item[\texttt{network\_payload}] Opera sobre eventos \texttt{network\_payload} 
    del \texttt{NetworkMonitor}. Aplica un diccionario de campo→regex 
    sobre el campo \texttt{data} del evento, que contiene el texto 
    ASCII imprimible extraído del payload del paquete. El matching se 
    realiza con \texttt{re.IGNORECASE}, de modo que el educador no 
    necesita contemplar variaciones de capitalización en los patrones.
    
    \item[\texttt{network\_connection}] Opera sobre eventos 
    \texttt{network\_connection} del \texttt{NetworkMonitor}. Detecta 
    conexiones TCP establecidas (paquete SYN-ACK) hacia un puerto de 
    destino específico, con filtro opcional por IP de destino.
    
    \item[\texttt{and}, \texttt{or}, \texttt{not}] Nodos lógicos 
    recursivos que permiten componer condiciones complejas sobre 
    cualquier combinación de targets atómicos. El campo 
    \texttt{condition: any|all} del escenario actúa como azúcar 
    sintáctica para listas planas de targets sin nodos lógicos.
\end{description}

Los tipos \texttt{http\_response\_contains} y 
\texttt{database\_query\_result} están definidos en el schema Pydantic 
pero no tienen evaluador implementado. Quedan documentados como trabajo 
futuro; el segundo en particular fue reemplazado funcionalmente por 
\texttt{network\_payload} con filtro \texttt{capture\_filter: "port 3306"}, 
que captura las queries directamente del tráfico TCP sin requerir acceso 
a la base de datos.

\section{Capa 4: Administración}

% [Descripción de la CLI con Typer + Rich, comando demo privesc,
% MonitorOrchestrator como punto de integración]